\chapter{Phao Thi Đại Chiến}
\chapterintro{Bi kịch của tiểu tự}{Chương này không ca ngợi quay cóp. Nó chỉ ghi lại đúng cái không khí run tay, viết chữ li ti, giấu phao rất kỹ rồi vào phòng thi lại quên sạch cách dùng.}

\poem{Bài 46}{Bài 46 - Chữ nhỏ hơn kiến bò}{Tiny cheat-sheet handwriting is its own absurd performance.}{Phao thi viết đêm trước\\Chữ li ti nhỏ hơn rận\\Đến khi xong tờ cuối cùng\\Mắt tôi mỏi gần phát khóc}

\poem{Bài 47}{Bài 47 - Giấu ở tay áo rất khéo}{A sleeve can hide many things, except panic.}{Phao thi giấu tay áo\\Tưởng đâu kín đáo lắm rồi\\Mới thấy giám thị đi lại\\Tim tôi nhảy trước chân tôi}

\poem{Bài 48}{Bài 48 - Vào phòng quên luôn ký hiệu}{A cheat sheet is useless when fear erases your own code.}{Phao thi soạn rất kỹ\\Mũi tên, ký hiệu đủ trò\\Vào phòng run tay một cái\\Chính mình đọc cũng không ra}

\poem{Bài 49}{Bài 49 - Giám thị nhìn là hóa đá}{The scariest thing in school is not the question sheet but the watchful eyes.}{Phao thi chưa kịp mở\\Đã nghe ghế kéo kèn kẹt\\Ngẩng lên gặp ngay mắt bác\\Tôi ngoan như tượng lập tức}

\poem{Bài 50}{Bài 50 - Xé phao như tuyết giữa sân}{Sometimes surrender is the only noble ending.}{Hết giờ ra ngoài ngõ\\Tôi xé phao nhỏ như mưa\\Gió bay lả tả một đoạn\\Thấy mình dại cũng vừa vừa}
